\chapter{Pruebas Unitarias}
\label{ch:testing}

\section{Herramientas utilizadas}

Las pruebas automatizadas del proyecto utilizan:

\begin{itemize}
    \item \textbf{JUnit 5} como framework de ejecución.
    \item \textbf{Mockito} para simular (\textit{mock}) repositorios y clientes Feign
    en las pruebas de la capa de servicio, aislándola de la base de datos real y de
    los demás microservicios.
    \item \textbf{@WebMvcTest} + \textbf{@MockitoBean} de Spring Boot Test para las
    pruebas de la capa de controlador, que verifican el código HTTP real devuelto por
    cada endpoint (\texttt{200}, \texttt{201}, \texttt{400}, \texttt{404}) sin levantar
    el contexto completo de Spring ni requerir base de datos.
    \item \textbf{JaCoCo} (\texttt{jacoco-maven-plugin}), integrado en el POM raíz, para
    medir la cobertura real de código ejercitada por las pruebas.
\end{itemize}

\section{Estrategia de pruebas}

Cada servicio cuenta con dos niveles de prueba:

\begin{enumerate}
    \item \textbf{Pruebas de servicio} (\texttt{XServiceImplTest}): verifican la lógica
    de negocio con los repositorios y clientes Feign mockeados, siguiendo el patrón
    \textit{Arrange-Act-Assert}. Cubren tanto el camino feliz como los casos de error
    (recurso no encontrado, argumento inválido).
    \item \textbf{Pruebas de controlador} (\texttt{XControllerTest}): verifican que
    cada endpoint devuelva el código HTTP correcto y la forma esperada de respuesta
    (incluyendo los enlaces HATEOAS), con el servicio mockeado mediante
    \texttt{@MockitoBean}.
\end{enumerate}

El siguiente ejemplo, tomado de las pruebas del microservicio de envíos, ilustra ambos
elementos: la verificación del camino feliz y la del registro correcto en el historial
de estados (la relación JPA agregada en la sección~\ref{sec:evolucion}):

\begin{lstlisting}[language=Java, caption={EnvioServiceImplTest.java -- fragmento}]
@Test
void deberiaCrearEnvioExitosamente() {
    Envio envio = new Envio();
    envio.setPedidoId(1L);
    when(pedidoClient.obtenerPedidoPorId(1L)).thenReturn(new PedidoRentDTO());
    when(envioRepository.save(any(Envio.class)))
        .thenAnswer(invocacion -> invocacion.getArgument(0));

    Envio resultado = envioService.crearEnvio(envio);

    assertEquals("PREPARACION", resultado.getEstado());
    verify(envioRepository, times(1)).save(any(Envio.class));
    verify(historialEstadoRepository, times(1)).save(any(HistorialEstado.class));
}
\end{lstlisting}

\section{Resultados de cobertura}

La ejecución de \texttt{mvn test} sobre los doce módulos del proyecto arroja
\textbf{168 pruebas} en total, todas exitosas. La Tabla~\ref{tab:cobertura} presenta la
cobertura de líneas medida con JaCoCo, calculada de dos formas distintas:

\begin{itemize}
    \item \textbf{Cobertura de lógica de negocio}: solo paquetes
    \texttt{controller}, \texttt{service} y \texttt{mapper} -- el código que realmente
    contiene reglas de negocio y flujo de la aplicación.
    \item \textbf{Cobertura del módulo completo}: incluye además clases de
    configuración (\texttt{WebConfig}, \texttt{OpenApiConfig}), la clase principal de
    arranque de Spring Boot, y los \texttt{DataSeeder} de Datafaker.
\end{itemize}

\begin{table}[h]
\centering
\begin{tabular}{lrr}
\toprule
\textbf{Microservicio} & \textbf{Lógica de negocio} & \textbf{Módulo completo} \\
\midrule
Catálogo & 100,0\% & 71,9\% \\
Clientes & 100,0\% & 70,6\% \\
Inventario & 100,0\% & 53,3\% \\
Notificaciones & 100,0\% & 70,3\% \\
Sucursales & 100,0\% & 70,1\% \\
Pedidos & 95,9\% & 39,6\% \\
Envíos & 93,8\% & 57,8\% \\
Pagos & 88,9\% & 49,5\% \\
Carrito & 81,0\% & 45,4\% \\
Reseñas & 84,3\% & 46,5\% \\
\midrule
\textbf{Total} & \textbf{93,6\%} & \textbf{55,2\%} \\
\bottomrule
\end{tabular}
\caption{Cobertura de líneas por microservicio (JaCoCo)}
\label{tab:cobertura}
\end{table}

La diferencia entre ambas columnas es intencional y se explica completamente: los
\texttt{DataSeeder} agregados para cumplir el requisito de Datafaker (sección
\ref{sec:evolucion}) son clases \texttt{CommandLineRunner} con lógica de reintento
contra Eureka, verificadas mediante ejecución real en Docker contra el stack completo
levantado -- no mediante pruebas unitarias con mocks, ya que su propósito es
precisamente ejercitar la comunicación real entre servicios al arrancar. Contarlas
como "no cubiertas" en un reporte de JaCoCo orientado a pruebas unitarias no refleja
que su comportamiento no haya sido validado, sino que se validó con una herramienta
distinta (verificación en vivo) y no cuenta como cobertura sobre las clases de
configuración (\texttt{WebConfig}, \texttt{OpenApiConfig}) tampoco aporta información
relevante, dado que no contienen lógica de negocio a probar. Por eso se reporta la
cobertura de lógica de negocio como el indicador principal de calidad de las pruebas.
